\chapter{Conclusion}

\paragraph{Summary of findings}
The findings in this thesis tie to the research questions posed in the introduction. Approximate Bayesian inference can be improved by geometry-aware modeling if the metric is chosen \emph{carefully}. We considered metrics constructed from the target distribution.
Paper~I concluded that the Riemannian Laplace approximation with the Fisher information matches strongly curved targets most accurately in small-dimensional problems.
Paper~II found that the correct choice of metric inside geodesic slice sampling captures curvature. If the target is multimodal, sampling using the inverse metrics improves movement between the modes.

Simplex geometry can be connected to Euclidean representations through stochastic interpolations and simplex-to-Euclidean bijections. Paper~III found that this makes Euclidean generative models applicable to discrete data, while allowing exact recovery of the original category with the $\arg\max$ operation.
Paper~IV showed that the Aitchison geometry together with the latent stochastic interpolation converts the problem into Gaussian process regression, which enables exact latent inference and gives accurate and calibrated predictions.

Overall, the thesis incorporated the geometry of the statistical problem into the method itself, which led to improved accuracy (while making it clear when these gains come at a computational cost) or to a latent simplification of the original problem. Next we give some of the limitations of the methods presented and avenues for future work.

\paragraph{Limitations and Outlook}
The main limitation of Papers I and II is computational complexity: numerical integration is expensive, and each integration step becomes cubic in complexity for the Fisher information metric. Future work can study (i) approximations of the geodesics, for example Taylor approximations \citep{hartmann:2023}, and (ii) amortization of the geodesics, for example geometry-aware variational approximations \citep{chevallier2022exponential}.

A limitation of Paper II is the lack of a guarantee of mode coverage. A more detailed theoretical study could show under which conditions the sampler can generate geodesics that reach other modes.

The main limitation of Paper III is that the method does not scale as well as discrete diffusion with respect to $K$. Recent work has shown that continuous diffusion does scale as well as discrete diffusion with additional improvements \citep{roos2026categorical, lee2026one}. Future work can explore whether training the flow directly with the negative categorical log likelihood (cross entropy) can improve our method in the same way. Another distinction with discrete diffusion is that some of these methods model the score only for a subset of categories \citep{Lou2024}, whereas we model all categories at the same time with the velocity field. It might be that restricting the velocity field to move between a subset of categories or to lower-dimensional projections improves the scalability of our method.

Paper IV shares the limitations of Gaussian process regression: it scales purely with the training data even with sparse approximations. Since our method is compatible with advances in Gaussian process regression, future work can test the scalability of the method on large classification tasks under scalable Gaussian process regression methods \citep{maronas2023efficient}.

The simplex-to-Euclidean bijections and the Aitchison geometry remain largely underutilized in modern machine learning. Other lines of future work include the sampling of densities constrained by a linear system. \citet{Diederen2025} construct a general bijection from convex polytopes to Euclidean space (which uses the ILR map). This can be used as a reparametrization of the target distribution defined on a convex set. In a similar way, high-dimensional categorical distributions can be converted to Euclidean target distributions with our procedure, which can ease sampling \citep{mohanty2026slithering}.

Classification with the softmax function leads to neural collapse \citep{kothapalli2023neural}. Future work can assess whether the ILR map avoids this behavior. This would open up the possibility of computing geodesics in the ILR-latent space, which might lead to improvements in the study of latent geometric representations \citep{yu2025connecting}.
